
\vspace{-0.5cm}

\begin{figure}[!h]
    \centering
    \vspace{-.5em}
    \includegraphics[width=.95\linewidth]{imgs/magentic_example.pdf}
    \vspace{-.5em}
    \caption{An illustration of the \team mutli-agent team completing a complex task from the GAIA benchmark. \team's Orchestrator agent creates a plan, delegates tasks to other agents, and tracks progress towards the goal, dynamically revising the plan as needed. The Orchestrator can delegate tasks to a FileSurfer agent to read and handle files, a WebSurfer agent to operate a web browser, or a Coder or Computer Terminal agent to write or execute code, respectively.}
    \label{fig:magentic-example}
\end{figure} 
\begin{abstract}

\noindent
Modern AI agents, driven by advances in large foundation models, promise to enhance our productivity and transform our lives by augmenting our knowledge and capabilities. To achieve this vision, AI agents must effectively plan, perform multi-step reasoning and actions, respond to novel observations, and recover from errors, to successfully complete complex tasks across a wide range of scenarios. In this work, we introduce {\em \team}, a high-performing open-source agentic system for solving such tasks. \team uses a multi-agent architecture where a lead agent, the \emph{Orchestrator}, plans, tracks progress, and re-plans to recover from errors. Throughout task execution, the Orchestrator also directs other specialized agents to perform tasks as needed, such as operating a web browser, navigating local files, or writing and executing Python code. Our experiments show that \team achieves statistically competitive performance to the state-of-the-art on three diverse and challenging agentic benchmarks: GAIA, AssistantBench, and WebArena. Notably, \team achieves these results without modification to core agent capabilities or to how they collaborate, demonstrating progress towards the vision of \emph{generalist agentic systems}. Moreover, \team's modular design allows agents to be added or removed from the team without additional prompt tuning or training, easing development and making it extensible to future scenarios. We provide an open-source implementation of \team, and we include \agbench{}, a standalone tool for agentic evaluation. \agbench{} provides built-in controls for repetition and isolation to run agentic benchmarks in a rigorous and contained manner -- which is important when agents' actions have side-effects. \team, \agbench{} and detailed empirical performance evaluations of \team,  including ablations and error analysis are available at \url{https://aka.ms/magentic-one}.
\end{abstract}
